%%%%%%%%%%  scan idx 196  |  folio 183  %%%%%%%%%%
Quitte à tronquer, on peut même supposer que $E^i=0$ pour $i<m$. D'après (I 4.16) appliqué pour $\mathcal C=\Mod$, $\mathcal C_0=\Lib$, on a alors $E^m\in\Ob\Loclib(S)$, d'où l'assertion.

c) L'hypothèse (1.1 (i))(resp. (1.1 (ii)respée)) étant vérifiée pour $\mathcal C_S=\Coh(S)$, $\mathcal C_0=\Lib$ d'après (2.4.3 a)) (resp. (2.1.1 b)) et le théorème A), on sait d'après (1.2 c) et d)) que le foncteur
\[
\mathrm K^b(\Loclib(S))\longrightarrow\mathrm D^b(\Coh(S))
\]
est pleinement fidèle. Donc, d'après (2.4.4), il en est de même du foncteur
\[
\mathrm K^b(\Loclib(S))\longrightarrow\mathrm D(S),
\]
dont l'image essentielle est $\mathrm D(S)_{\parf}$ en vertu de b). Par un raisonnement analogue à celui fait dans la preuve de (2.4.3 c)), on voit que le foncteur $\mathrm K^b(\Loclib(S))\to\mathrm D(S)$ induit par $\Gamma_S$ est pleinement fidèle et a pour image essentielle $\mathrm D(S)_{\parf}$, et l'on achève la démonstration comme en a).

\textbf{Remarque 2.4.6.} Le rédacteur ignore si dans (2.4.4) on peut remplacer $\mathrm D^b$ par $\mathrm D^{-}$, et par suite si (2.4.5 a)) est encore vrai lorsqu'on remplace $\mathrm K^{-,b}$ par $\mathrm K^{-}$ et $\mathrm D^b$ par $\mathrm D^{-}$.

Voici, pour terminer, un complément intéressant à (2.4.5). Donnons d'abord une définition :

\textbf{Définition 2.4.7.} On appelle \emph{espace pseudo-analytique complexe} tout espace annelé localement isomorphe à un espace annelé du type $(Y,\mathcal O_X|Y)$, où $X$ est un espace analytique complexe, et $Y$ une partie fermée de $X$.

Le faisceau structural d'un espace pseudo-analytique complexe est donc cohérent.

On dit qu'un espace pseudo-analytique complexe est de Stein si le topos correspondant l'est (2.4.1).



% Exposé II continuation macros
\providecommand{\Qcoh}{\operatorname{Qcoh}}
\providecommand{\Mod}{\operatorname{Mod}}
\providecommand{\Coh}{\operatorname{Coh}}
\renewcommand{\Ob}{\operatorname{ob}}
\providecommand{\Hom}{\operatorname{Hom}}
\providecommand{\Ext}{\operatorname{Ext}}
\providecommand{\Spec}{\operatorname{Spec}}
\providecommand{\qcoh}{\mathrm{qcoh}}
\providecommand{\parf}{\mathrm{parf}}
\providecommand{\coh}{\mathrm{coh}}
\providecommand{\Loclib}{\operatorname{Loclib}}
\providecommand{\Lib}{\operatorname{Lib}}
\providecommand{\im}{\operatorname{Im}}

